% 敏感性分析表
\documentclass{article}
\usepackage[utf8]{inputenc}
\usepackage{booktabs}
\usepackage{amsmath}
\usepackage{caption}
\captionsetup[table]{labelformat=empty}
\renewcommand{\arraystretch}{1.2}

\begin{document}

\begin{table}[htbp]
\centering
\caption{Sensitivity Analysis: Effect of Meditation-state Label Mapping Value on Model Performance (15-fold LOSO)}
\label{tab:sensitivity}
\begin{tabular}{lccc}
\toprule
\textbf{Meditation Value} & \textbf{MSE} & \textbf{Pearson $r$} & \textbf{$\Delta$ MSE vs.\ 0.6} \\
\midrule
0.5 & $0.1571 \pm 0.1209$ & $0.675 \pm 0.298$ & $+0.0047$ \\
0.6 (Default) & $0.1524 \pm 0.1395$ & $0.723 \pm 0.287$ & --- \\
0.7 & $0.1372 \pm 0.1187$ & $0.716 \pm 0.299$ & $-0.0152$ \\
\bottomrule
\end{tabular}
\vspace{0.5em}

\footnotesize{
\textit{Note:} All three settings yield comparable performance (MSE range: 0.137--0.157), confirming that the framework is robust to the specific choice of the meditation mapping value within the physiological range. The default value of 0.6 is retained following the empirical soft-label encoding logic described in Section~\ref{sec:dataset_task}.
}
\end{table}
\end{document}
